\documentclass[a4paper]{report}
% Shared preamble for the English probability and statistics documents.
% Based on the original course preamble. Compile with pdfLaTeX.
\usepackage[top=1.5cm,bottom=2cm,left=2cm,right=2cm]{geometry}
\usepackage[T1]{fontenc}
\usepackage{lmodern}
\usepackage[utf8]{inputenc}
\usepackage[english]{babel}

\usepackage{amsmath}
\usepackage{amssymb}
\usepackage{amsthm}
\usepackage{mathrsfs}
\usepackage{mathtools}
\usepackage{etoolbox}
\usepackage{tikz}
\usetikzlibrary{shapes.multipart,arrows.meta}
\usepackage{xcolor}
\usepackage{enumerate}
\usepackage[shortlabels]{enumitem}
\usepackage{hyperref}
\usepackage{version}
\usepackage{bbold}
\usepackage{fontawesome5}
\usepackage{fancyhdr}
\lfoot{}
\usepackage{natbib}
\usepackage{graphicx}
\usepackage{pdfpages}
\usepackage{authblk}
    
\usepackage{blindtext}
\usepackage{multicol}
\usepackage{collcell}
\usepackage{wrapfig}
\usepackage{blkarray}
\usepackage{cancel}
\usepackage{appendix}
\usepackage{float}
\makeatletter
\input{tex/assets/slashbox.sty}
\makeatother
\hypersetup{unicode,colorlinks,citecolor=red,filecolor=black,
  linkcolor=red,urlcolor=cyan}
\setlength{\columnsep}{1cm}
\pagenumbering{gobble}

\def\numtext#1{%
  \ifcase#1\or one\or two\or three\or four\or five\or six\or seven\or eight\or nine\or ten\or
  eleven\or twelve\or thirteen\or fourteen\or fifteen\or sixteen\or seventeen\or eighteen\or nineteen\or twenty\or Lots\fi}

\newcommand{\Byte}[1]{%
  \def\nodecontents{}%
  \foreach[count=\l] \k in {#1} {\xappto{\nodecontents}{\noexpand\nodepart{\numtext{\l}}\k}}%
  \node[rectangle split,rectangle split parts=8,rectangle split horizontal,draw] at (2,2){\nodecontents};}

\newcommand{\D}{{\mathrm{d}}}
\newcommand{\cE}{{\mathcal{E}}}
\newcommand{\cG}{{\mathcal{G}}}
\newcommand{\cX}{{\mathcal{X}}}
\newcommand{\cY}{{\mathcal{Y}}}
\newcommand{\cB}{{\mathcal{B}}}
\newcommand{\cF}{{\mathcal{F}}}
\newcommand{\cO}{{\mathcal{O}}}
\newcommand{\cS}{{\mathcal{S}}}
\newcommand{\cT}{{\mathcal{T}}}
\newcommand{\given}{\, \big| \,}
\newcommand{\cL}{{\mathscr{L}}}
\newcommand{\cD}{{\mathcal{D}}}
\newcommand{\cA}{{\mathscr{A}}}
\newcommand{\cP}{{\mathscr{P}}}
\newcommand{\cR}{{\mathscr{R}}}
\newcommand{\A}{{\mathcal A}}
\newcommand{\one}{\mathbb{1}}
\newcommand{\mathbbm}[1]{\mathbb{#1}}

\providecommand{\C}{}
\renewcommand{\C}{\mathbb C}
\newcommand{\R}{\mathbb R}
\newcommand{\Q}{\mathbb Q}
\newcommand{\N}{\mathbb N}
\newcommand{\T}{\mathbb T}
\newcommand{\F}{\mathbb F}
\renewcommand{\P}{\mathbb P}
\newcommand{\Z}{\mathbb Z}
\newcommand{\E}{\mathbb {E}}
\newcommand{\Pbb}[1]{\P\left( #1\right)}
\newcommand{\Ebb}[1]{\E\left[ #1\right]}

\newcommand{\lb}{\left (}
\newcommand{\rb}{\right )}
\newcommand{\lbb}{\left [}
\newcommand{\rbb}{\right ]}
\newcommand{\labs}{\left |}
\newcommand{\rabs}{\right |}
\newcommand{\lbrb}[1]{\lb #1 \rb}
\newcommand{\lbbrbb}[1]{\lbb#1\rbb}
\newcommand{\labsrabs}[1]{\labs#1\rabs}
\newcommand{\lbbrb}[1]{\lbb#1\rb}
\newcommand{\lbrbb}[1]{\lb#1\rbb}
\newcommand{\langlerangle}[1]{\langle#1\rangle}
\newcommand{\lbcurly}{\{}
\newcommand{\rbcurly}{\}}
\newcommand{\curly}[1]{\lbcurly#1\rbcurly}
\newcommand{\abs}{\labsrabs}
\newcommand{\ind}[1]{1_{#1}}
\newcommand{\indep}{\perp \!\!\! \perp}
\renewcommand{\d}{\, \mathrm{d}}

\newtheorem{theorem}{Theorem}
\newtheorem{problemCnt}{Lemma}
\theoremstyle{definition}
\newtheorem{problem}[problemCnt]{Problem}
\newtheorem{problembook}[problemCnt]{\faBook \,  Problem}
\newtheorem{problempencil}[problemCnt]{\faPencil* \,  Problem}
\newtheorem{problemstar}[problemCnt]{\faStar\,  Problem}
\newtheorem{problemstarpencil}[problemCnt]{\faStar \, \faPencil* \,   Problem}
\newtheorem{probleminteresting}[problemCnt]{\faSmile[regular] \,   Problem}
\newtheorem{solution}{Solution}

\newcommand{\fwcell}[1]{\makebox[\arraystretch\normalbaselineskip]{$#1$}}
\newenvironment{enumeratei}{\begin{enumerate}[(i)]}{\end{enumerate}}
\newenvironment{enumeratea}{\begin{enumerate}[(a)]}{\end{enumerate}}
\newenvironment{enumeraten}{\begin{enumerate}[\upshape 1.]}{\end{enumerate}}

% Compatibility helpers needed by a few files in the combined archive.
\AtBeginDocument{\noindent{\footnotesize English translation prepared with ChatGPT (OpenAI, GPT-6), 26 October 2026, for translation of the Bulgarian text into English and LaTeX formatting.
The Bulgarian originals remain the reference; translation may contain errors.}\par\smallskip}
\newcommand{\MissingGraphic}[1]{%
  \begin{center}\fbox{\parbox{0.82\linewidth}{\centering Missing source image: \texttt{#1}}}\end{center}}



\usepackage{tcolorbox}
\setcounter{section}{0}
\renewcommand{\thesection}{\arabic{section}}
\begin{document}
\noindent Probability Theory and Mathematical Statistics, 2024/2025 \hfill Homework 4

\hrulefill

\begin{it}
The maximum total across all homework assignments this semester is 50 points,
corresponding to a bonus of 0.5 on the tutorial grade.


This assignment is based on an entrance examination for
the Centrale-Sup\'elec group of French universities. See
\url{https://www.concours-centrale-supelec.fr/sujets-rapports}, under Math\'ematiques
1 for 2018. It is therefore distributed under the same licence,
CC BY-NC-SA 2.0: \url{https://creativecommons.org/licenses/by-nc-sa/2.0/}.

\end{it}

Notes:
\begin{itemize}
\item Throughout, \(N\), \(k\), and \(d\) are integers at least 2,
and $t$ is a real number.
\item For positive integers \(m,n\), \(\mathcal M_{m,n}(\mathbb R)\) denotes the real \(m\times n\) matrices.
  \item 
For brevity,
we use
  $[1,N]$,
to mean
  $\{1,\dots,N\}$.
\item For a matrix \(A\), \(A^T\) denotes its transpose.
\item All random objects are defined on the finite probability space \((\Omega,\mathcal A,\mathbb P)\).
\item For a Euclidean space \(E\), such as $\mathbb R^n$, write \(\langle\cdot,\cdot\rangle\) and \(\|\cdot\|\) for its inner product and induced norm.
\end{itemize}

\section*{Problem statement and objective}

What is the largest number of points in $\mathbb R^n$
that can all be at equal distances from one another?
In the plane ($n=2$) the answer is 3, forming an equilateral triangle;
for $n=3$ it is 4, forming a regular tetrahedron. One can check that in $\mathbb R^n$ it is $n+1$.

Thus, for example, 50 vectors in $\mathbb R^{100}$ can be mutually equidistant. The preceding observation shows that depicting their positions in two dimensions is impossible
without changing the distances! This is a relatively simple manifestation of the \emph{curse of dimensionality}, a collective term for the often counterintuitive properties of high-dimensional spaces.

This assignment asks:
\begin{tcolorbox}[colframe=black, colback=white, boxrule=0.5mm, sharp corners]
Can $N$ points in a high-dimensional Euclidean space be linearly projected into a lower-dimensional space without substantially distorting their pairwise distances?
\end{tcolorbox}

More precisely, let \(v_1,\ldots,v_N\in\mathbb R^d\) be distinct. For \(\varepsilon\in(0,1)\), call a linear map \(f:\mathbb R^d\to\mathbb R^k\) an \emph{\(\varepsilon\)-isometry} for these vectors if
\[
\forall (i,j)\in[ 1,N] ^2:\quad
(1-\varepsilon)\,\|v_i-v_j\|\le \|f(v_i)-f(v_j)\|\le (1+\varepsilon)\,\|v_i-v_j\|.
\]

\medskip

\textbf{Objective:} For which \(k\) does there exist a linear \(\varepsilon\)-isometry \(f:\mathbb R^d\to\mathbb R^k\) for \(v_1,\ldots,v_N\)?

\medskip

We will prove the following theorem, first established by William B. Johnson and Joram Lindenstrauss in 1984:
\begin{tcolorbox}[colframe=black, colback=white, boxrule=0.5mm, sharp corners]
\emph{There is an absolute constant \(c>0\) such that, for all integers \(N,d\geq2\) and every set of \(N\) distinct vectors in \(\mathbb R^d\), the condition}
\[
k\ge \frac{c\ln(N)}{\varepsilon^2},
\]
\emph{is sufficient for an \(\varepsilon\)-isometry \(f:\mathbb R^d\to\mathbb R^k\) for \(v_1,\ldots,v_N\) to exist.}
\end{tcolorbox}

The only known proofs of this theorem use probabilistic methods. Section \ref{sec:1} develops preliminary results on convexity and probability. Section \ref{sec:2} proves a concentration result of Talagrand, which Section \ref{sec:3} applies to the Johnson--Lindenstrauss theorem.
\pagebreak

\section{Preliminary results}\label{sec:1}

\subsection{Projection onto a closed convex set}
Let \(E\) be a Euclidean space.
\begin{enumerate}
\item Let \(a,b\in E\). Prove and geometrically interpret
  \[
  \|a+b\|^2+\|a-b\|^2 = 2\lbrb{\|a\|^2 + \|b\|^2}.
  \]
\item Deduce that if \(u,v,v'\in E\), \(v\ne v'\), and \(\|u-v\|=\|u-v'\|\), then
  \[
  \Bigl\|u-\frac{v+v'}{2}\Bigr\| < \|u-v\|.
  \]
\item Let \(F\subset E\) be nonempty and closed, and \(u\in E\). Prove that there exists \(v\in F\) such that
  \[
  \forall w\in F:\quad \|u-v\|\le \|u-w\|.
  \]
\item Deduce that if \(C\subset E\) is nonempty, closed and convex, then for every \(u\in E\) there is a unique \(v\in C\) such that
  \[
  \forall w\in C:\quad \|u-v\|\le \|u-w\|.
  \]
Call \(v\) the \emph{projection} of \(u\) onto \(C\), and write
  \(\,d(u,C):=\|u-v\|\).
\end{enumerate}

\subsection{H\"older's inequality for expectations}
Let \(p,q\) be positive real numbers satisfying
\[
\frac{1}{p} + \frac{1}{q} = 1.
\]
\begin{enumerate}
\item Use concavity of the logarithm to prove, for positive \(a,b\),
  \[
  ab \le \frac{a^p}{p} + \frac{b^q}{q}.
  \]
\item Deduce that random variables \(X,Y\) on a finite probability space satisfy
  \[
  \Ebb{|XY|}\le \lbrb{\Ebb{|X|^p}}^{1/p}\,
  \lbrb{\Ebb{|Y|^q}}^{1/q}.
  \]
(\emph{First consider $\Ebb{|X|^p}=\Ebb{|Y|^q}=1$.})
\end{enumerate}


\subsection{Random variables with sub-Gaussian tails}
Let \(X:\Omega\to\mathbb R\) and suppose there are \(a,b>0\) such that, for all \(t>0\),
\[
\mathbb{P}(|X|\ge t)\le a\exp(-bt^2).
\]
\begin{enumerate}
\item Prove that
  \[
  \Ebb{X^2}
  = 2\int_{0}^{+\infty} t\,\mathbb{P}(|X|\ge t)\,\D t.
  \]
\item Show that \(\Ebb{X^2}\leq a/b\).
\item For \(\delta\) with \(\delta^2\leq a/b\), prove
  \[
  \mathbb{P}(|X+\delta|\ge t)\le \mathbb{P}(|X|\ge t-|\delta|).
  \]
\item Prove that
  \[
  -b(t-|\delta|)^2 \le a - \frac12bt^2.
  \]
\item Deduce that for \(t\geq|\delta|\),
  \[
  \mathbb{P}(|X+\delta|\ge t)
  \le a\exp(a)\exp\Bigl(-\frac12bt^2\Bigr).
  \]
\item Show that this inequality also holds for \(0\leq t<|\delta|\).
\end{enumerate}

\bigskip

\section{Talagrand's inequality}
\label{sec:2}
Let \(E\) be an \(n\)-dimensional Euclidean space, \(n\geq1\), with orthonormal basis \((e_1,\ldots,e_n)\). Let \(\varepsilon_1,\ldots,\varepsilon_n\) be independent Rademacher variables, meaning
$\P(\epsilon_1=1)=\P(\epsilon_1=-1)=1/2$, and set
\[
X:=\sum_{i=1}^n \varepsilon_i\,e_i.
\]
Our aim is to show that for every nonempty closed convex \(C\subset E\),
\begin{equation}
    \label{eq: Talagrand}
\mathbb{P}\lbrb{X\in C}\,
\Ebb{\exp\lbrb{\frac18\,d(X,C)^2}}\le1,
\end{equation}
and use this to establish Talagrand's inequality \eqref{eq:Talagrand2}.

\subsection{Two special cases}
\begin{enumerate}
\item Consider the case where \(C\) does not intersect \[X(\Omega):=\{x\in E:X(\omega)=x
\text{ for some }\omega\in\Omega\},\] that is, $C\cap X(\Omega)=\emptyset$.
  
For the remaining parts of this subsection, assume \(C\cap X(\Omega)\) consists of a single vector \(u\).
\item Show that \(\frac14d(X,u)^2\sim\mathrm{Bin}(n,1/2)\).
\item Find the expectation of \(\exp(\frac18d(X,u)^2)\) and show that it is at most \(2^n\).
\item Prove \(d(X,C)\leq d(X,u)\), and establish
  \eqref{eq: Talagrand}.
\end{enumerate}

\subsection{Induction}
Now suppose \(C\subset E\) is closed and convex, with at least two elements in \(C\cap X(\Omega)\).
Without loss of generality, assume two of those vectors differ in their last coordinate.
We prove \eqref{eq: Talagrand} by induction on \(n=\dim E\).
\begin{enumerate}
\item Consider \(n=1\).
\end{enumerate}

\subsection{Properties of \(C_{+1}\) and \(C_{-1}\)}
Let
\[
E':=\mathrm{span}(e_1,\dots,e_{n-1}),\quad
\pi:E\to E'\text{ be the orthogonal projection},
\quad
X':=\pi(X)=\sum_{i=1}^{n-1}\varepsilon_i\,e_i.
\]
For \(t\in\{-1,1\}\), let
\[
H_t := E' + t\,e_n,\quad
C_t := \pi\bigl(C\cap H_t\bigr).
\]
\begin{enumerate}
\item Show that for every \(x'\in E'\) and \(t\in\{-1,1\}\),
  \[
  x'\in C_t \iff x' + t\,e_n\in C.
  \]
\item Prove that \(C_{+1},C_{-1}\) are nonempty closed convex subsets of \(E'\).
\item For \(t\in\{-1,1\}\), let \(Y_t\) be the projection of \(X'\) onto \(C_t\).
Show that
  \[
  \mathbb{P}(X\in C)
  = \frac12\,\mathbb{P}(X'\in C_{+1})
  + \frac12\,\mathbb{P}(X'\in C_{-1}).
  \]
\end{enumerate}

\subsection{A key inequality}
Let \(0\leq\lambda\leq1\).
\begin{enumerate}
\item Show that
  \[
  d(X,C)\le \Bigl\|\,(1-\lambda)\bigl(Y_{\varepsilon_n}+\varepsilon_n e_n\bigr)
  +\lambda\bigl(Y_{-\varepsilon_n}-\varepsilon_n e_n\bigr)-X\Bigr\|.
  \]
\item Then prove that
  \[
  d(X,C)^2\le 4\lambda^2+\Bigl\|(1-\lambda)(Y_{\varepsilon_n}-X')
  +\lambda(Y_{-\varepsilon_n}-X')\Bigr\|^2
  \]
and
  \[
  d(X,C)^2\le 4\lambda^2+(1-\lambda)\|Y_{\varepsilon_n}-X'\|^2
  +\lambda\|Y_{-\varepsilon_n}-X'\|^2,
  \]
that is,
  \[
  d(X,C)^2\le 4\lambda^2+(1-\lambda)d(X',C_{\varepsilon_n})^2
  +\lambda\,d(X',C_{-\varepsilon_n})^2.
  \]
\end{enumerate}

\subsection{Conditional expectations}
Let
\[
p_+ := \mathbb{P}(X'\in C_{+1}),\quad
p_- := \mathbb{P}(X'\in C_{-1}),
\]
and assume without loss of generality that \(p_+\geq p_-\).
\begin{enumerate}
\item Show that \(p_->0\).
\item Prove that for every \(\lambda\in[0,1]\),
  \[
  \begin{aligned}
    &\Ebb{\exp\lbrb{\frac18\,d(X,C)^2}\mid \varepsilon_n=-1}\\
    &\quad\le
    \exp\lbrb{\frac{\lambda^2}{2}}
    \;\Ebb{%
      \lbrb{\exp\lbrb{\frac18\,d(X',C_{-1})^2}}^{1-\lambda}
      \;\lbrb{\exp\lbrb{\frac18\,d(X',C_{+1})^2}}^{\lambda}
    }.
  \end{aligned}
  \]
\item Deduce that
  \[
    \Ebb{\exp\lbrb{\frac18\,d(X,C)^2}\mid \varepsilon_n=-1}
    \le
    \exp\lbrb{\frac{\lambda^2}{2}}
    \;\Ebb{\exp\lbrb{\frac18\,d(X',C_{-1})^2}}^{1-\lambda}
    \;\Ebb{\exp\lbrb{\frac18\,d(X',C_{+1})^2}}^{\lambda}.
  \]
\item Using the induction hypothesis, justify
  \[
    \Ebb{\exp\lbrb{\frac18\,d(X,C)^2}\mid \varepsilon_n=1}
    \le \frac{1}{p_+}.
  \]
and then show that, for every \(\lambda\in[0,1]\),
  \[
    \Ebb{\exp\lbrb{\frac18\,d(X,C)^2}}
    \le
    \frac12\;\lbrb{
      \frac1{p_+}
      + \exp\lbrb{\frac{\lambda^2}{2}}
        \;\frac1{(p_-)^{1-\lambda}}
        \;\frac1{(p_+)^{\lambda}}
    }.
  \]
\end{enumerate}

\subsection{Optimization}
\begin{enumerate}
\item Let \(\lambda=1-p_-/p_+\). Show that
  \[
    \Ebb{\exp\lbrb{\frac18\,d(X,C)^2}}
    \le
    \frac{1}{2p_+}\;\lbrb{
      1 + \exp\lbrb{\frac{\lambda^2}{2}}\,(1-\lambda)^{\lambda-1}
    }.
  \]
\item Prove that for every \(x\in[0,1)\),
  \[
    \frac{x^2}{2} + (x-1)\,\ln\lbrb{1-x}
    \;\le\;
    \ln\lbrb{2+x} - \ln\lbrb{2-x}.
  \]
\item Deduce that for every \(x\in[0,1)\),
  \[
    1 + \exp\lbrb{\frac{x^2}{2}}\,(1-x)^{x-1}
    \;\le\;
    \frac{4}{2-x}.
  \]
\item Complete the proof of \eqref{eq: Talagrand}.
\item Prove Talagrand's inequality:
    \begin{tcolorbox}[colframe=black, colback=white, boxrule=0.5mm, sharp corners]
\emph{For every nonempty closed convex \(C\subset E\) and every \(t>0\),
    \begin{equation}\label{eq:Talagrand2}
      \mathbb{P}(X\in C)\,\mathbb{P}(d(X,C)\ge t)
      \;\le\;
      \exp\lbrb{-\frac{t^2}{8}}.
    \end{equation}}
    \end{tcolorbox}
\end{enumerate}


\bigskip

\section{The Johnson--Lindenstrauss theorem}\label{sec:3}
In this section consider
\[
  E = \mathcal{M}_{k,d}(\mathbb{R}),
\]
equipped with the inner product
\(\langle A, B\rangle_F := \mathrm{tr}(A^T B)\)
and its induced (Frobenius) norm \(\|\cdot\|_F\). Identify \(\mathbb R^d\) with \(\mathcal M_{d,1}(\mathbb R)\), viewing vectors as column matrices. Let \(u\in\mathbb R^d\), \(\|u\|=1\), and define
\[
  g(M) := \|M\,u\|,
  \quad
  M\in\mathcal{M}_{k,d}(\mathbb{R}).
\]
Let
\[
  X = (\varepsilon_{ij})_{1\le i\le k,\;1\le j\le d}
\]
be a random matrix in \(\mathcal M_{k,d}(\mathbb R)\) whose entries \(\varepsilon_{ij}\) are independent Rademacher variables.

\subsection{A concentration inequality}
\begin{enumerate}
\item Show that
  \[
    C = \{M \mid g(M)\le r\}
  \]
is a closed convex subset of \(\mathcal M_{k,d}(\mathbb R)\).
\item Prove that every matrix \(M\) satisfies
  \[
    \|M\,u\| \le \|M\|_F.
  \]
\item Let \(r,t\in\mathbb R\), \(t>0\). Show that \(d(M,C)<t\) implies \(g(M)<r+t\).
\item Deduce that
  \[
    \mathbb{P}(g(X)\le r)\,\mathbb{P}(g(X)\ge r+t)
    \;\le\;
    \exp\lbrb{-\frac{t^2}{8}}.
  \]
\end{enumerate}

\subsection{Medians}
\begin{enumerate}
\item Explain why \(g(X)\) has at least one median.
\item Prove that for every \(t>0\),
  \[
    \mathbb{P}\lbrb{|g(X)-m|\ge t}
    \;\le\;
    4\,\exp\lbrb{-\frac{t^2}{8}},
  \]
and that
  \[
    \Ebb{(g(X)-m)^2} \le 32.
  \]
\item Show that \(\Ebb{g(X)^2}=k\) and \(\Ebb{g(X)}\leq\sqrt k\).
\item Deduce that
  \[
    (\sqrt{k}-m)^2 \le \Ebb{(g(X)-m)^2}.
  \]
\end{enumerate}

\subsection{A key lemma}
\begin{enumerate}
\item Show that for every \(t>0\),
  \[
    \mathbb{P}\lbrb{|g(X)-\sqrt{k}|\ge t}
    \;\le\;
    4\,\exp(4)\,\exp\lbrb{-\frac{t^2}{16}}.
  \]
\item Define \(A_k=X/\sqrt k\), and let \(\varepsilon\in(0,1)\), \(\delta\in(0,1/2)\), with
  \[
    k \ge \frac{160\,\ln(1/\delta)}{\varepsilon^2}.
  \]
Prove that for every vector \(u\)
of unit norm,
  \[
    \mathbb{P}\lbrb{\big|\,\|A_k\,u\|-1\big| > \varepsilon} < \delta.
  \]
\end{enumerate}

\subsection{Conclusion}
Let $v_1,\ldots,v_N$ be distinct vectors in $\mathbb R^d$. For \((i,j)\in[1,N]^2\), \(i<j\), define
\[
  E_{ij}
  := \lbrb{%
    (1-\varepsilon)\,\|v_i-v_j\|\le\|A_kv_i-A_kv_j\|\le(1+\varepsilon)\,\|v_i-v_j\|%
  }.
\]
\begin{enumerate}
\item Show that \(\mathbb P(E_{ij}^c)<\delta\), where \(E_{ij}^c=\Omega\setminus E_{ij}\).
\item Deduce that
  \[
\mathbb{P}\lbrb{\bigcap_{i<j}E_{ij}}
    \;\ge\;
    1 - \frac{N(N-1)}{2}\,\delta.
  \]
\item Prove the Johnson--Lindenstrauss theorem.
\end{enumerate}

\end{document}

